%=============================================================================
%  CURSO DE SQL Y BASES DE DATOS RELACIONALES
%  Presentación Beamer — Edición Profesional y Optimizada
%  Autor: Miyoshi Mikami — CNBV
%=============================================================================

\documentclass[12pt, aspectratio=169, UTF8]{beamer}

% ─── Tema base ────────────────────────────────────────────
\usetheme{Madrid}

% ─── Paleta de colores institucional y tecnológica ────────
\definecolor{sqlblue}{RGB}{15, 52, 96}       % Azul marino profundo

\definecolor{sqlgray}{RGB}{44, 62, 80}       % Gris oscuro para texto
\definecolor{sqlyellow}{RGB}{241, 196, 15}   % Amarillo acento

\definecolor{sqlred}{RGB}{192, 57, 43}       % Rojo alerta
\definecolor{codefg}{RGB}{228, 232, 237}     % Texto código (mejor contraste)
\definecolor{codebg}{RGB}{13, 17, 23}        % Fondo código (GitHub Dark real)

% ─── Aplicar paleta al tema Beamer ────────────────────────
\setbeamercolor{palette primary}  {bg=sqlblue,   fg=white}
\setbeamercolor{palette secondary}{bg=sqlblue!85,fg=white}
\setbeamercolor{palette tertiary} {bg=sqlblue!70,fg=white}
\setbeamercolor{palette quaternary}{bg=sqlblue,  fg=white}
\setbeamercolor{structure}        {fg=sqlblue}
\setbeamercolor{frametitle}       {bg=sqlblue, fg=white}
\setbeamercolor{title}            {fg=white}
\setbeamercolor{subtitle}         {fg=sqllight}
\setbeamercolor{section in toc}   {fg=sqlblue}

% Bloques estructurados estándar
\setbeamercolor{block title}      {bg=sqlaccent, fg=white}
\setbeamercolor{block body}       {bg=sqllight,  fg=sqlgray}
\setbeamercolor{block title alerted}{bg=sqlred,  fg=white}
\setbeamercolor{block body alerted} {bg=sqlred!8, fg=sqlgray}
\setbeamercolor{block title example}{bg=sqlgreen,fg=white}
\setbeamercolor{block body example} {bg=sqlgreen!8,fg=sqlgray}

% ─── Fuentes y Paquetes Esenciales ────────────────────────
\usepackage[T1]{fontenc}
\usepackage[spanish, es-tabla]{babel}
\usepackage{lmodern}
\usepackage{microtype}

% Íconos y símbolos matemáticos
\usepackage{amssymb, amsmath}
\usepackage{graphicx}
\usepackage{booktabs}
\usepackage{tabularx}
\usepackage{array}

% NOTA: Se eliminó 'enumitem' porque rompe el comportamiento nativo 
% de las listas en Beamer (animaciones, espaciados y templates).

% ─── Código SQL con Resaltado Inteligente ─────────────────
\usepackage{listings}

\lstdefinestyle{sqlstyle}{
  language         = SQL,
  basicstyle       = \ttfamily\footnotesize\color{codefg},
  keywordstyle     = \bfseries\color{sqlaccent},
  stringstyle      = \color{sqlyellow},
  commentstyle     = \itshape\color{gray!65},
  numberstyle      = \tiny\color{gray!50},
  numbers          = left,
  numbersep        = 6pt,
  stepnumber       = 1,
  backgroundcolor  = \color{codebg},
  frame            = single,
  frameround       = tttt,
  framexleftmargin = 2pt,
  rulecolor        = \color{sqlblue!30},
  tabsize          = 4,
  showstringspaces = false,
  breaklines       = true,
  breakatwhitespace= true,
  xleftmargin      = 0.4cm,
  xrightmargin     = 0.1cm,
  aboveskip        = 6pt,
  belowskip        = 6pt,
  extendedchars    = true,
  literate         = {á}{{\'a}}1 {é}{{\'e}}1 {í}{{\'i}}1 {ó}{{\'o}}1 {ú}{{\'u}}1 {ñ}{{\~n}}1 {Ú}{{\'U}}1 {Ñ}{{\~N}}1 {¿}{{{\textquestiondown}}}1 {ó}{{\'o}}1,
  morekeywords     = {
    DATABASE, DATABASES, TABLES, SHOW, USE, ENGINE, CHARSET,
    AUTO_INCREMENT, UNSIGNED, TINYINT, SMALLINT, MEDIUMINT,
    BIGINT, DECIMAL, FLOAT, DOUBLE, DATETIME, TIMESTAMP,
    VARCHAR, TEXT, BLOB, ENUM, SET, DESCRIBE, DESC,
    FOREIGN KEY, REFERENCES, ON DELETE, ON UPDATE,
    CASCADE, SET NULL, RESTRICT, TRIGGER, VIEW,
    PROCEDURE, CALL, DELIMITER, SAVEPOINT, ROLLBACK,
    BEGIN, COMMIT, START, TRANSACTION, SIGNAL, SQLSTATE,
    MATCH, AGAINST, FULLTEXT, SPATIAL, UNIQUE, INDEX,
    LAST_INSERT_ID, COALESCE, IFNULL, FORMAT, CONCAT,
    CONCAT_WS, YEAR, DATEDIFF, ROUND, UPPER, LOWER,
    LENGTH, TRIM, SUBSTRING, NOW, BETWEEN, LIKE, IN,
    HAVING, GROUP, ORDER, LIMIT, OFFSET, UNION, ALL,
    INNER, LEFT, RIGHT, JOIN, ON, AS, DISTINCT, CASE,
    WHEN, THEN, ELSE, END, COUNT, SUM, AVG, MIN, MAX,
    INSERT, UPDATE, DELETE, CREATE, DROP, ALTER, TRUNCATE,
    ADD, RENAME, MODIFY, COLUMN, TABLE, PRIMARY, KEY,
    NOT, NULL, DEFAULT, CHECK, WITH, ROLLUP, IF, EXISTS, CURDATE
  }
}
\lstset{style=sqlstyle}

% ─── Contenedores Estilizados (tcolorbox) ─────────────────
\usepackage[many]{tcolorbox}
\tcbuselibrary{skins}

\newtcolorbox{sqlbox}[1]{
  enhanced,
  colback   = codebg,
  colframe  = sqlaccent,
  coltitle  = white,
  fonttitle = \bfseries\small\ttfamily,
  title     = #1,
  top=3pt, bottom=3pt, left=5pt, right=5pt,
  sharp corners,
  boxrule=1pt,
  shadow={2pt}{-2pt}{0pt}{black!20}
}

% ─── Personalización Estética de Listas ───────────────────
\setbeamertemplate{itemize items}[circle]
\setbeamercolor{itemize item}{fg=sqlaccent}
\setbeamercolor{itemize subitem}{fg=sqlyellow}
\setbeamertemplate{enumerate items}[default]
\setbeamercolor{enumerate item}{fg=sqlaccent}

% ─── Footer y Limpieza de Barra de Navegación ─────────────
\setbeamertemplate{navigation symbols}{}
\setbeamertemplate{footline}{
  \leavevmode%
  \hbox{%
  \begin{beamercolorbox}[wd=.333333\paperwidth,ht=2.25ex,dp=1ex,left,leftskip=10pt]{palette primary}%
    \usebeamerfont{author in head/foot}\insertshortauthor\ (\insertshortinstitute)
  \end{beamercolorbox}%
  \begin{beamercolorbox}[wd=.333333\paperwidth,ht=2.25ex,dp=1ex,center]{palette secondary}%
    \usebeamerfont{title in head/foot}\insertshorttitle
  \end{beamercolorbox}%
  \begin{beamercolorbox}[wd=.333333\paperwidth,ht=2.25ex,dp=1ex,right,rightskip=10pt]{palette primary}%
    \usebeamerfont{date in head/foot}\insertshortdate\hspace*{2em}
    \insertframenumber{} / \inserttotalframenumber
  \end{beamercolorbox}%
  }\vskip0pt%
}

% ─── Automatización de Páginas de Transición ──────────────
\AtBeginSection[]{
  \begin{frame}
    \vfill
    \centering
    \begin{beamercolorbox}[sep=12pt,center,shadow=true,rounded=true]{palette primary}
      \usebeamerfont{title}\Huge\insertsectionhead\par%
    \end{beamercolorbox}
    \vfill
  \end{frame}
}

% ─── Metadata del Documento ───────────────────────────────
\title[Curso de SQL]{Curso de SQL y Bases de Datos Relacionales}
\subtitle{Inserción, Consulta y Administración Eficiente de Información}
\author[Miyoshi Mikami]{Miyoshi Mikami \\ \small \href{mailto:miyoshikamium@gmail.mx}{\texttt{miyoshikamium@gmail.mx}}}
\institute[CNBV]{Comisión Nacional Bancaria y de Valores}
\date[Junio 2026]{Versión 4.5 — Junio 2026}

%=============================================================================
\begin{document}
%=============================================================================

% ── PORTADA ──────────────────────────────────────────────
{
\setbeamercolor{background canvas}{bg=sqlblue}
\begin{frame}[plain]
  \titlepage
\end{frame}
}

% ── ÍNDICE ───────────────────────────────────────────────
\begin{frame}{Contenido del Curso}
  \tableofcontents[hideallsubsections]
\end{frame}

%=============================================================================
\section{Historia e Introducción}
%=============================================================================

\begin{frame}{El Problema de los Datos}
  \begin{columns}[T]
    \column{0.55\textwidth}
      \textbf{Evolución del almacenamiento humano:}
      \begin{itemize}
        \item Tablillas de arcilla sumerias ($\sim$3000 a.C.)
        \item Registros contables del Imperio Romano
        \item Libros de contabilidad medievales
      \end{itemize}
      \vspace{6pt}
      El siglo XX transformó esto en una \textbf{disciplina formal}:
      \begin{itemize}
        \item Crecimiento exponencial de registros masivos
        \item Necesidad crítica de velocidad y fiabilidad
      \end{itemize}
    \column{0.42\textwidth}
      \begin{block}{Pregunta clave}
        ¿Cómo organizar grandes volúmenes de datos para que múltiples programas accedan de forma concurrente, eficiente y segura?
      \end{block}
  \end{columns}
\end{frame}

\begin{frame}{Modelos Pre-Relacionales (1950–1970)}
  \begin{columns}[T]
    \column{0.48\textwidth}
      \begin{block}{Modelo Jerárquico}
        \begin{itemize}
          \item Estructura rígida de \textbf{árbol}
          \item IBM IMS (1966) — Misión Apolo
          \item Relaciones N:M imposibles sin duplicidad redundante
        \end{itemize}
      \end{block}
    \column{0.48\textwidth}
      \begin{block}{Modelo de Red (CODASYL 1969)}
        \begin{itemize}
          \item Punteros físicos entre registros
          \item Soporta relaciones N:M nativas
          \item Complejidad de mantenimiento extrema
        \end{itemize}
      \end{block}
  \end{columns}
  \vspace{6pt}
  \begin{alertblock}{Problema fundamental: Dependencia físico-lógica}
    Los programas requerían conocer la ubicación exacta en el disco. Cualquier modificación estructural rompía la compatibilidad de la aplicación.
  \end{alertblock}
\end{frame}

\begin{frame}{Edgar F. Codd y el Modelo Relacional (1970)}
  \begin{columns}[T]
    \column{0.56\textwidth}
      En junio de \textbf{1970}, Codd publica su artículo histórico en la \emph{ACM}:
      \vspace{4pt}
      \begin{block}{Modelo Relacional}
        Representación uniforme mediante \textbf{tablas bidimensionales}. Operaciones basadas en el \textbf{álgebra relacional}, abstrayendo la arquitectura física subyacente.
      \end{block}
      \vspace{4pt}
      \textbf{Pilares de la revolución:}
      \begin{enumerate}
        \item Independencia total físico-lógica
        \item Rigor matemático (Teoría de conjuntos)
        \item Enfoque netamente \emph{declarativo}
      \end{enumerate}
    \column{0.40\textwidth}
      \begin{exampleblock}{Premio Turing 1981}
        Otorgado a Edgar F. Codd por cambiar de forma permanente las ciencias de la computación.
      \end{exampleblock}
      \begin{block}{Las 12 Reglas (1985)}
        Métricas obligatorias para evaluar la madurez relacional de un SGBD.
      \end{block}
  \end{columns}
\end{frame}

\begin{frame}{Nacimiento de SQL}
  \begin{columns}[T]
    \column{0.56\textwidth}
      \begin{itemize}
        \item \textbf{1973} — IBM desarrolla el proyecto \emph{System R}
        \item Chamberlin y Boyce definen \textbf{SEQUEL} (\emph{Structured English Query Language})
        \item Rediseñado a \textbf{SQL} por derechos de marca
        \item \textbf{1979} — Oracle V2: Lanzamiento del primer SGBD comercial del mercado
        \item \textbf{1983} — Comercialización formal de IBM DB2
        \item \textbf{1986} — Adopción del primer estándar oficial ANSI/ISO
      \end{itemize}
    \column{0.40\textwidth}
      \begin{alertblock}{Pronunciación Oficial}
        \begin{itemize}
          \item Coloquial anglosajón: \textit{"sequel"}
          \item Estándar formal ISO: \textbf{"ese-cu-ele"}
        \end{itemize}
        Adoptamos la nomenclatura formal y estandarizada.
      \end{alertblock}
  \end{columns}
\end{frame}

\begin{frame}{Estándares del Lenguaje SQL}
  \begin{center}
  \small
  \renewcommand{\arraystretch}{1.15}
  \begin{tabularx}{0.95\textwidth}{@{} l l X @{}}
    \toprule
    \textbf{Versión} & \textbf{Año} & \textbf{Aportaciones e Innovaciones Clave} \\
    \midrule
    SQL-86    & 1986 & Formalización inicial de las sintaxis DML y DDL básico. \\
    SQL-92    & 1992 & Introducción formal de \texttt{JOIN} explícitos y subconsultas. \\
    SQL:1999  & 1999 & Soporte de recursión (\texttt{WITH}), \textit{triggers} y POO básica. \\
    SQL:2003  & 2003 & Funciones analíticas de ventana (\texttt{OVER}) e integración XML. \\
    SQL:2016  & 2016 & Compatibilidad e interoperabilidad nativa con formato JSON. \\
    SQL:2023  & 2023 & Integración con bases de datos de grafos y uso extendido de \textit{arrays}. \\
    \bottomrule
  \end{tabularx}
  \end{center}
  \begin{alertblock}{Portabilidad}
    Ningún motor comercial cubre el estándar al 100\%; las funciones avanzadas requieren adecuación según la distribución (PostgreSQL, MySQL, Oracle).
  \end{alertblock}
\end{frame}

%=============================================================================
\section{Fundamentos de Bases de Datos}
%=============================================================================

\begin{frame}{¿Qué es una Base de Datos?}
  \begin{block}{Definición Formal}
    Colección \textbf{organizada, persistente y compartible} de registros interrelacionados, diseñada metodológicamente para \textbf{eliminar redundancias} y preservar consistencia.
  \end{block}
  \vspace{4pt}
  \begin{columns}[T]
    \column{0.48\textwidth}
      \textbf{Dimensiones estructurales:}
      \begin{itemize}
        \item \textbf{Persistencia:} Resguardo robusto en almacenamiento estable.
        \item \textbf{Concurrencia:} Mecanismos de acceso seguro multiusuario.
        \item \textbf{Normalización:} Estructuras lógicas eficientes y limpias.
      \end{itemize}
    \column{0.48\textwidth}
      \begin{exampleblock}{Caso de Aplicación Práctica}
        \textbf{Sin Base de Datos:} Redundancia descontrolada en archivos aislados.\\[4pt]
        \textbf{Modelo Relacional:} Identidades únicas mapeadas mediante relaciones atómicas.
      \end{exampleblock}
  \end{columns}
\end{frame}

\begin{frame}{Sistema de Gestión de Bases de Datos (SGBD)}
  \begin{block}{Definición Técnica}
    Software sofisticado que actúa como capa de abstracción intermedia entre el almacenamiento de bajo nivel y las solicitudes lógicas de los usuarios.
  \end{block}
  \begin{columns}[T]
    \column{0.48\textwidth}
      \textbf{Modelo de Tres Niveles ANSI/SPARC:}
      \begin{enumerate}
        \item \textbf{Interno / Físico:} Organización binaria en disco.
        \item \textbf{Conceptual:} Tablas, esquemas, restricciones.
        \item \textbf{Externo / Vistas:} Interfaces lógicas personalizadas.
      \end{enumerate}
    \column{0.48\textwidth}
      \textbf{Funcionalidad Core:}
      \begin{itemize}
        \item Motor de optimización de ejecución.
        \item Aislamiento integral concurrente.
        \item Recuperación automática ante fallos de energía.
      \end{itemize}
  \end{columns}
\end{frame}

\begin{frame}{Equivalencia Léxica y Claves}
  \begin{center}
  \small
  \renewcommand{\arraystretch}{1.2}
  \begin{tabularx}{0.9\textwidth}{@{} X X X @{}}
    \toprule
    \textbf{Término Académico} & \textbf{Terminología SQL} & \textbf{Concepto Visual} \\
    \midrule
    Relación  & Tabla            & Matriz de datos \\
    Tupla     & Fila / Registro  & Línea horizontal de registros \\
    Atributo  & Columna / Campo  & Componente vertical estructurado \\
    Dominio   & Tipo de Dato     & Especificación y reglas de valor \\
    Esquema   & Definición       & Metadatos estructurales de la tabla \\
    \bottomrule
  \end{tabularx}
  \end{center}
  \begin{columns}[T]
    \column{0.48\textwidth}
      \begin{block}{Clave Primaria (PK)}
        Atributo que identifica de manera única e inequívoca una fila. Prohíbe nulos (\texttt{NOT NULL}).
      \end{block}
    \column{0.48\textwidth}
      \begin{block}{Clave Foránea (FK)}
        Vínculo referencial apuntando a una PK externa. Asegura consistencia relacional.
      \end{block}
  \end{columns}
\end{frame}

\begin{frame}{El Paradigma de las Propiedades ACID}
  \begin{columns}[T]
    \column{0.48\textwidth}
      \begin{block}{\textbf{A}tomicidad}
        La transacción se procesa de forma binaria: se ejecuta al 100\% o se revierte por completo.
      \end{block}
      \begin{block}{\textbf{C}onsistencia}
        Los cambios mueven el sistema únicamente entre estados lógicos válidos del esquema.
      \end{block}
    \column{0.48\textwidth}
      \begin{block}{\textbf{I}solamiento}
        Las transacciones concurrentes se ejecutan sin interferir en sus estados intermedios.
      \end{block}
      \begin{block}{\textbf{D}urabilidad}
        Una vez confirmado con \texttt{COMMIT}, el registro es imborrable ante fallos críticos.
      \end{block}
  \end{columns}
  \begin{exampleblock}{Ejemplo Financiero Basal}
    Una transferencia bancaria exige el descuento en Origen \textbf{y} abono en Destino. Si falla la red a la mitad, la \textbf{atomicidad} ejecuta un rollback inmediato.
  \end{exampleblock}
\end{frame}

\begin{frame}{Sublenguajes Especializados en SQL}
  \begin{center}
  \small
  \renewcommand{\arraystretch}{1.25}
  \begin{tabularx}{0.92\textwidth}{@{} l l X @{}}
    \toprule
    \textbf{Sublenguaje} & \textbf{Sigla} & \textbf{Propósito y Comandos Cardinales} \\
    \midrule
    Data Definition Language   & \textbf{DDL} & Define y altera estructuras de datos: \texttt{CREATE, ALTER, DROP}. \\
    Data Manipulation Language & \textbf{DML} & Operación de datos residentes: \texttt{SELECT, INSERT, UPDATE, DELETE}. \\
    Data Control Language      & \textbf{DCL} & Gestión de esquemas de seguridad y accesos: \texttt{GRANT, REVOKE}. \\
    Transaction Control Lang.  & \textbf{TCL} & Orquestación de bloques operativos: \texttt{COMMIT, ROLLBACK, SAVEPOINT}. \\
    \bottomrule
  \end{tabularx}
  \end{center}
  \begin{alertblock}{Naturaleza Declarativa}
    El analista programa el \textbf{QUÉ} requiere. El Optimizador del motor calcula analíticamente la ruta algorítmica (\textbf{CÓMO}) más rápida.
  \end{alertblock}
\end{frame}

%=============================================================================
\section{Consulta de Datos: SELECT}
%=============================================================================

\begin{frame}[fragile]{Arquitectura de una Consulta SELECT}
  \begin{sqlbox}{Sintaxis Canónica}
SELECT [ALL | DISTINCT] lista_columnas
FROM tabla [[AS] alias]
[WHERE condiciones_de_filtrado_lineal]
[GROUP BY expresiones_de_agrupamiento]
[HAVING filtros_de_agregacion_grupal]
[ORDER BY criterio_orden [ASC | DESC]]
[LIMIT registros_maximos OFFSET desfasamiento];
  \end{sqlbox}
  \vspace{2pt}
  \textbf{Orden de Ejecución Lógica del Motor Interno:}
  \begin{center}
    \footnotesize
    \color{sqlblue}\textbf{FROM} $\rightarrow$ \textbf{WHERE} $\rightarrow$ \textbf{GROUP BY} $\rightarrow$ \textbf{HAVING} $\rightarrow$ \textbf{SELECT} $\rightarrow$ \textbf{DISTINCT} $\rightarrow$ \textbf{ORDER BY} $\rightarrow$ \textbf{LIMIT}
  \end{center}
\end{frame}

\begin{frame}[fragile]{Proyección Eficiente y Campos Calculados}
  \begin{sqlbox}{Casos de Proyección Atómica}
-- Proyección explícita limpia
SELECT nombre, departamento, salario 
FROM empleados;

-- Inclusión de expresiones algebraicas dinámicas
SELECT nombre,
       salario * 12 AS salario_anual,
       salario * 0.10 AS provision_bono,
       'Vigente' AS estatus_operativo
FROM empleados;
  \end{sqlbox}
  \begin{alertblock}{Restricción de Buenas Prácticas}
    Se debe omitir el uso de \texttt{SELECT *} en arquitecturas productivas. Incrementa la latencia de red, deshabilita optimizaciones de índices e induce errores al alterar esquemas.
  \end{alertblock}
\end{frame}

\begin{frame}[fragile]{Filtros Lineales mediante WHERE}
  \begin{columns}[T]
    \column{0.40\textwidth}
      \small
      \renewcommand{\arraystretch}{1.2}
      \begin{tabular}{@{} c l @{}}
        \toprule
        \textbf{Operador} & \textbf{Significado Analítico} \\
        \midrule
        \texttt{=}   & Equivalencia estricta \\
        \texttt{<>} o \texttt{!=} & Discrepancia lógica \\
        \texttt{>}, \texttt{<}   & Mayor / Menor estricto \\
        \texttt{>=}, \texttt{<=} & Límites inclusivos \\
        \bottomrule
      \end{tabular}
      \vspace{6pt}
      \begin{exampleblock}{Mecánica}
        Evalúa registro por registro; solo los bloques verdaderos (\texttt{TRUE}) se integran a la memoria temporal.
      \end{exampleblock}
    \column{0.56\textwidth}
      \begin{sqlbox}{Uso de Predicados de Filtrado}
SELECT nombre, cargo, salario
FROM empleados
WHERE departamento = 'TI'
  AND salario > 45000;
      \end{sqlbox}
  \end{columns}
\end{frame}

\begin{frame}[fragile]{Ordenamiento y Paginación de Datos}
  \begin{sqlbox}{Sintaxis de Ordenamiento y Control de Bloques}
-- Ordenamiento jerárquico multicriterio
SELECT nombre, departamento, salario
FROM empleados
ORDER BY departamento ASC, salario DESC;

-- Paginación de datos eficiente (Página 2, tamaño 3)
SELECT nombre, departamento
FROM empleados
ORDER BY nombre ASC
LIMIT 3 OFFSET 3;
  \end{sqlbox}
  \begin{alertblock}{Determinismo en Paginación}
    Utilizar \texttt{LIMIT} sin definir un \texttt{ORDER BY} explícito quita el determinismo a la consulta. El motor puede regresar conjuntos inconsistentes según el estado interno de las páginas del disco.
  \end{alertblock}
\end{frame}

%=============================================================================
\section{Modificadores: Parte 1}
%=============================================================================

\begin{frame}[fragile]{Filtros Avanzados y Operadores de Patrón}
  \begin{sqlbox}{Cláusula DISTINCT para Valores Únicos}
SELECT DISTINCT departamento FROM empleados;
  \end{sqlbox}
  \vspace{4pt}
  \begin{sqlbox}{Evaluación de Cadenas mediante LIKE y Comodines}
SELECT nombre FROM empleados 
WHERE nombre LIKE 'A%';     -- Cadenas con inicio en 'A'
WHERE nombre LIKE '%ez';    -- Apellidos con terminación 'ez'
WHERE nombre LIKE '%ar%';   -- Subcadenas internas
WHERE nombre LIKE '_o%';    -- Fila con 'o' fijada en segunda posición
  \end{sqlbox}
  \begin{center}
    \footnotesize
    \texttt{\%} $\rightarrow$ Comodín para n caracteres intermedios. \quad \texttt{\_} $\rightarrow$ Comodín de espacio posicional único.
  \end{center}
\end{frame}

\begin{frame}[fragile]{Precedencia y Operadores Booleanos}
  \begin{columns}[T]
    \column{0.40\textwidth}
      \small
      \renewcommand{\arraystretch}{1.2}
      \begin{tabular}{@{} l l @{}}
        \toprule
        \textbf{Operador} & \textbf{Jerarquía de Evaluación} \\
        \midrule
        \texttt{NOT} & Prioridad Máxima \\
        \texttt{AND} & Prioridad Intermedia \\
        \texttt{OR}  & Prioridad Baja \\
        \bottomrule
      \end{tabular}
      \vspace{10pt}
      \begin{alertblock}{Peligro Lógico}
        Omitir el aislamiento lógico de las sentencias altera drásticamente el resultado debido a la prioridad intrínseca del \texttt{AND}.
      \end{alertblock}
    \column{0.56\textwidth}
      \begin{sqlbox}{Agrupamiento Explicito con Paréntesis}
SELECT nombre, salario
FROM empleados
WHERE (departamento = 'TI' AND cargo = 'Gerente')
   OR (departamento = 'Ventas' AND salario > 30000);
      \end{sqlbox}
  \end{columns}
\end{frame}

%=============================================================================
\section{Modificadores: Parte 2}
%=============================================================================

\begin{frame}[fragile]{Funciones de Agregación y Métricas de Grupo}
  \begin{sqlbox}{Métricas de Agregación Estándar}
SELECT departamento,
       COUNT(*) AS total_nomina,
       ROUND(AVG(salario), 2) AS salario_promedio,
       SUM(salario) AS costo_operativo,
       MIN(salario) AS sueldo_piso,
       MAX(salario) AS sueldo_techo
FROM empleados
GROUP BY departamento
HAVING AVG(salario) > 40000
ORDER BY costo_operativo DESC;
  \end{sqlbox}
  \begin{exampleblock}{Diferenciación Clave en Optimización}
    \textbf{WHERE:} Evalúa y descarta registros individuales \textit{antes} del proceso de consolidación en memoria.\\
    \textbf{HAVING:} Evalúa propiedades agregadas sobre los bloques estructurados \textit{después} del armado grupal.
  \end{exampleblock}
\end{frame}

\begin{frame}[fragile]{Rangos, Listas y Funciones de Formato}
  \begin{sqlbox}{Predicados de Inclusión: IN y BETWEEN}
-- Evaluación eficiente de conjuntos cerrados
SELECT nombre FROM empleados WHERE departamento IN ('Ventas', 'TI', 'Finanzas');

-- Evaluación inclusiva de límites vectoriales
SELECT nombre FROM empleados WHERE salario BETWEEN 35000 AND 52000;
  \end{sqlbox}
  \vspace{4pt}
  \begin{sqlbox}{Operaciones de Salida y Concatenación}
SELECT nombre AS empleado,
       CONCAT(nombre, ' — ', cargo) AS etiqueta_puesto
FROM empleados;
  \end{sqlbox}
\end{frame}

\begin{frame}[fragile]{Estructuras Condicionales y Tratamiento de Nulos}
  \begin{sqlbox}{Evaluación Lógica mediante CASE WHEN}
SELECT nombre, salario,
       CASE 
            WHEN salario >= 55000 THEN 'Estructura Senior'
            WHEN salario >= 45000 THEN 'Estructura Mid'
            ELSE 'Estructura Junior'
       END AS tabulador_puesto
FROM empleados;
  \end{sqlbox}
  \vspace{4pt}
  \begin{sqlbox}{Manejo de la Ausencia de Datos (Valores NULL)}
-- Sustitución directa parametrizada
SELECT IFNULL(telefono, 'Sin Registro') FROM empleados;

-- Evaluación de cortocircuito secuencial vectorizada
SELECT COALESCE(tel_celular, tel_fijo, email, 'Incomunicable') AS contacto_oficial
FROM empleados;
  \end{sqlbox}
\end{frame}

%=============================================================================
\section{Consulta de Datos Relacionados}
%=============================================================================

\begin{frame}{Álgebra de JOINs y Combinaciones Lógicas}
  \begin{center}
  \small
  \renewcommand{\arraystretch}{1.2}
  \begin{tabularx}{0.95\textwidth}{@{} l X l @{}}
    \toprule
    \textbf{Operación SQL} & \textbf{Comportamiento de Intersección} & \textbf{Aplicación Práctica} \\
    \midrule
    \texttt{INNER JOIN}  & Registros coincidentes en ambos universos. & Cruce estricto de entidades. \\
    \texttt{LEFT JOIN}   & Universo Izquierdo íntegro + coincidencias. & Auditoría e identificación de huérfanos. \\
    \texttt{RIGHT JOIN}  & Universo Derecho íntegro + coincidencias.  & Idéntico a Left invirtiendo el orden. \\
    \texttt{UNION}       & Fusión vertical sin duplicidad de filas.  & Consolidación de catálogos homogéneos. \\
    \texttt{UNION ALL}   & Fusión vertical masiva con duplicidad.    & Concatenación veloz sin costo de ordenamiento. \\
    \bottomrule
  \end{tabularx}
  \end{center}
\end{frame}

\begin{frame}[fragile]{Mecánica de Combinaciones en Código}
  \begin{sqlbox}{INNER JOIN — Coincidencias Exactas}
SELECT c.nombre, p.pedido_id, p.monto
FROM clientes AS c
INNER JOIN pedidos AS p ON c.cliente_id = p.cliente_id;
  \end{sqlbox}
  \vspace{4pt}
  \begin{sqlbox}{LEFT JOIN — Inclusión Total y Detección de Nulos}
SELECT c.nombre, p.pedido_id, p.monto
FROM clientes AS c
LEFT JOIN pedidos AS p ON c.cliente_id = p.cliente_id;

-- Detección analítica de clientes inactivos (sin transacciones)
-- WHERE p.pedido_id IS NULL;
  \end{sqlbox}
\end{frame}

%=============================================================================
\section{Escritura de Datos}
%=============================================================================

\begin{frame}[fragile]{Operaciones de Manipulación DML}
  \begin{sqlbox}{Escritura, Mutación y Eliminación Masiva}
-- Inserción multi-registro atómica
INSERT INTO productos (nombre, precio, stock) VALUES 
('Teclado Mecánico', 850.00, 40),
('Monitor 24"', 3200.00, 15);

-- Modificación con alcance delimitado
UPDATE productos 
SET precio = precio * 1.10 
WHERE categoria_id = 3;

-- Borrado físico selectivo
DELETE FROM productos 
WHERE stock = 0 AND categoria_id = 5;
  \end{sqlbox}
\end{frame}

\begin{frame}{Contraste Técnico: DELETE vs. TRUNCATE}
  \begin{center}
  \small
  \renewcommand{\arraystretch}{1.2}
  \begin{tabularx}{0.95\textwidth}{@{} l X X @{}}
    \toprule
    \textbf{Dimensión Técnica} & \textbf{DELETE} & \textbf{TRUNCATE} \\
    \midrule
    Clasificación Formal       & Lenguaje DML & Lenguaje DDL \\
    Capacidad de Filtrado      & Permite cláusula \texttt{WHERE}. & Remueve la tabla completa. \\
    Manejo Transaccional       & Registra logs (\texttt{ROLLBACK} viable). & Operación directa sobre bloques de disco. \\
    Desempeño e Impacto        & Lento (evaluación fila por fila). & Instantáneo (libera punteros de almacenamiento). \\
    Contadores Secuenciales    & Preserva el estado de \texttt{AUTO\_INC}. & Resetea contadores al origen base. \\
    Mecanismos de Disparo      & Ejecuta \textit{Triggers} asociados. & Ignora disparadores lógicos. \\
    \bottomrule
  \end{tabularx}
  \end{center}
  \begin{alertblock}{Directiva de Seguridad Monetaria}
    Usa \texttt{DELETE} dentro de bloques de transacciones controladas en entornos de producción. Reserva \texttt{TRUNCATE} para procesos aislados de carga masiva de datos (ETL).
  \end{alertblock}
\end{frame}

%=============================================================================
\section{Administración de Tablas}
%=============================================================================

\begin{frame}[fragile]{Inicialización de Esquemas e Internacionalización}
  \begin{sqlbox}{Inicialización del Contenedor de Base de Datos}
CREATE DATABASE IF NOT EXISTS tienda_online
    CHARACTER SET utf8mb4
    COLLATE utf8mb4_unicode_ci;

USE tienda_online;
  \end{sqlbox}
  \vspace{4pt}
  \begin{alertblock}{Importancia Crítica de utf8mb4}
    El juego de caracteres tradicional \texttt{utf8} de MySQL restringe el almacenamiento a 3 bytes por símbolo. Para asegurar compatibilidad con glifos internacionales modernos, ideogramas y codificación emoji, es obligatorio forzar el uso de \textbf{\texttt{utf8mb4}}.
  \end{alertblock}
\end{frame}

\begin{frame}{Cardinalidad y Restricciones de Integridad}
  \begin{columns}[T]
    \column{0.32\textwidth}
      \begin{block}{Relaciones 1:1}
        \footnotesize Entidad con extensión única. Mapeada mediante una Foreign Key restringida explícitamente con indexación de unicidad \texttt{UNIQUE}.
      \end{block}
    \column{0.32\textwidth}
      \begin{block}{Relaciones 1:N}
        \footnotesize Estructura clásica Padre-Hijo. La clave foránea se inyecta directamente dentro de la tabla dependiente.
      \end{block}
    \column{0.32\textwidth}
      \begin{block}{Relaciones N:M}
        \footnotesize Dependencia distributiva asociativa. Exige la creación imperativa de una \textbf{tabla intermedia} relacional.
      \end{block}
  \end{columns}
  \vspace{10pt}
  \begin{alertblock}{Regla Operativa en Sistemas Transaccionales}
    \textbf{Orden de Inserción:} Escribir primero las tablas base (Padres) antes de asociar elementos en las dependientes.\\
    \textbf{Orden de Eliminación:} Operar en sentido inverso para no violar las restricciones de integridad referencial.
  \end{alertblock}
\end{frame}

%=============================================================================
% \section{Conceptos Avanzados}
%=============================================================================

\begin{frame}[fragile]{Estrategias de Indexación (INDEX)}
  \begin{sqlbox}{Tipos de Índices y Optimización de Consultas}
-- Índice B-Tree estándar para búsquedas rápidas
CREATE INDEX idx_apellido ON clientes(apellido);

-- Índice compuesto ideal para filtros ordenados combinados
CREATE INDEX idx_cat_precio ON productos(categoria_id, precio);

-- Índice con restricción estricta de unicidad lógica
CREATE UNIQUE INDEX uq_correo ON usuarios(correo);

-- Índice Full-Text para optimización de motores de búsqueda semántica
CREATE FULLTEXT INDEX ft_desc ON productos(descripcion);
  \end{sqlbox}
\end{frame}

\begin{frame}[fragile]{Automatización Interna mediante Triggers}
  \begin{sqlbox}{Trigger Corporativo de Auditoría de Precios}
DELIMITER $$
CREATE TRIGGER trg_auditoria_precio
AFTER UPDATE ON productos
FOR EACH ROW
BEGIN
    IF OLD.precio <> NEW.precio THEN
        INSERT INTO auditoria_precios (producto_id, precio_ant, precio_new, usuario, fecha)
        VALUES (NEW.producto_id, OLD.precio, NEW.precio, USER(), NOW());
    END IF;
END$$
DELIMITER ;
  \end{sqlbox}
\end{frame}

\begin{frame}[fragile]{Abstracciones Lógicas: Vistas y Procedimientos}
  \begin{sqlbox}{Encapsulamiento Lógico Mediante VIEW}
CREATE VIEW vista_pedidos_activos AS
SELECT p.pedido_id, c.nombre AS cliente, p.total
FROM pedidos AS p
INNER JOIN clientes AS c ON p.cliente_id = c.cliente_id
WHERE p.estado = 'activo';
  \end{sqlbox}
  \vspace{4pt}
  \begin{sqlbox}{Invocación de Procedimientos Almacenados (Stored Procedures)}
-- Ejecución modular parametrizada de alta seguridad
CALL sp_resumen_cliente(7, @num_pedidos, @monto);
SELECT @num_pedidos AS pedidos, @monto AS total;
  \end{sqlbox}
\end{frame}

\begin{frame}[fragile]{Control Transaccional Complejo}
  \begin{sqlbox}{Uso Avanzado de Bloques Operativos y Savepoints}
START TRANSACTION;

INSERT INTO pedidos (cliente_id, total, fecha) VALUES (12, 1500.00, CURDATE());
SAVEPOINT sp_pedido_base;

UPDATE productos SET stock = stock - 3 WHERE producto_id = 5;

-- Si las condiciones del inventario son inválidas:
-- ROLLBACK TO sp_pedido_base;

COMMIT;
  \end{sqlbox}
  \begin{exampleblock}{Mecanismo de Aislamiento de Motores (InnoDB)}
    El nivel nativo \textbf{REPEATABLE READ} imposibilita lecturas sucias asegurando la inmutabilidad de los datos consultados a lo largo de toda la transacción activa.
  \end{exampleblock}
\end{frame}

%=============================================================================
\section{Conexión desde Código}
%=============================================================================

\begin{frame}[fragile]{Conectores de Aplicación e Inyección SQL}
  \begin{sqlbox}{Implementación Orientada a Objetos en Python}
import mysql.connector
conn = mysql.connector.connect(
    host="localhost", user="app_user", password="s3cr3to", database="tienda_online"
)
cursor = conn.cursor(dictionary=True)
cursor.execute("SELECT nombre, precio FROM productos LIMIT 5")
for fila in cursor.fetchall():
    print(f"Producto: {fila['nombre']} | Precio: {fila['precio']}")
  \end{sqlbox}
  \vspace{2pt}
  \begin{alertblock}{Mitigación Total de Vulnerabilidades OWASP (SQL Injection)}
\begin{lstlisting}[language=Python, numbers=none, frame=none]
# INSEGURO (Vulnerable a ataques destructivos por concatenación):
query = "SELECT * FROM usuarios WHERE user = '" + input_usuario + "'"

# SEGURO (Uso Mandatorio de Consultas Parametrizadas y Enlazadas):
cursor.execute("SELECT * FROM usuarios WHERE user = %s", (input_usuario,))
\end{lstlisting}
  \end{alertblock}
\end{frame}

%=============================================================================
\section{Próximos Pasos}
%=============================================================================

\begin{frame}{Estrategias de Cierre y Práctica Profesional}
  \begin{columns}[T]
    \column{0.48\textwidth}
      \begin{block}{Hitos de Consolidación Técnico}
        \begin{itemize}
          \item Diseñar diagramas entidad-relación (DER) normalizados en tercera forma normal (3FN).
          \item Practicar perfiles de ejecución analizando consultas complejas mediante el comando \texttt{EXPLAIN}.
        \end{itemize}
      \end{block}
    \column{0.48\textwidth}
      \begin{exampleblock}{Laboratorios de Práctica Sugeridos}
        \begin{itemize}
          \item Instalación local y configuración óptima de motores usando entornos aislados con Docker.
          \item Resolución de retos relacionales avanzados en plataformas interactivas especializadas.
        \end{itemize}
      \end{exampleblock}
  \end{columns}
  \vspace{10pt}
  \centering
  \textbf{\Large ¡Muchas gracias por su atención!} \\
  \vspace{5pt}
  \footnotesize Comisión Nacional Bancaria y de Valores — Coordinación de Ciencia de Datos
\end{frame}

\end{document}



